\documentclass[11pt]{article}

\usepackage[margin=0.9in]{geometry}
\usepackage{amsmath,amssymb}
\usepackage{booktabs}
\usepackage{enumitem}
\usepackage{listings}
\usepackage{xcolor}
\usepackage{hyperref}

\definecolor{backcolour}{rgb}{0.95,0.95,0.92}
\definecolor{codeblue}{rgb}{0.05,0.15,0.55}
\definecolor{codegreen}{rgb}{0.0,0.45,0.0}

\lstset{
    backgroundcolor=\color{backcolour},
    basicstyle=\ttfamily\small,
    keywordstyle=\color{codeblue}\bfseries,
    commentstyle=\color{codegreen},
    stringstyle=\color{purple},
    breaklines=true,
    frame=single,
    showstringspaces=false,
    tabsize=4
}

\title{CPSC 250 \\ Simple Linear Regression with StatsModels}
\author{Edward Brash}
\date{}

\begin{document}

\maketitle

\section{Motivation}

Previously, we used \texttt{scipy.optimize.curve\_fit()} for fitting data.

In this lecture, we introduce a different approach using:

\begin{itemize}
    \item \texttt{pandas}
    \item \texttt{statsmodels}
\end{itemize}

The goals are:

\begin{enumerate}
    \item obtain parameter uncertainties automatically,
    \item use a cleaner data structure,
    \item simplify regression workflows.
\end{enumerate}

\section{Why Pandas?}

A \texttt{pandas DataFrame} is similar to a spreadsheet.

Each column has a label and stores related data.

This is often easier to work with than raw NumPy matrices.

\section{Creating a DataFrame}

We first import pandas:

\begin{lstlisting}[language=Python]
import pandas as pd
\end{lstlisting}

Then create a DataFrame:

\begin{lstlisting}[language=Python]
df = pd.DataFrame({

    "x": [5, 10, 15, 20],
    "y": [5, 12, 20, 27]

})
\end{lstlisting}

The keys:

\begin{lstlisting}[language=Python]
"x"
"y"
\end{lstlisting}

become the column labels.

\section{Accessing Columns}

We can print the entire table:

\begin{lstlisting}[language=Python]
print(df)
\end{lstlisting}

Or access a specific column:

\begin{lstlisting}[language=Python]
print(df["x"])
\end{lstlisting}

This returns only the x-values.

\section{Preparing for Linear Regression}

Suppose we want a linear model:

\[
y = \beta_0 + \beta_1 x
\]

where:

\begin{itemize}
    \item \(\beta_0\) is the intercept,
    \item \(\beta_1\) is the slope.
\end{itemize}

\section{Importing StatsModels}

We typically import:

\begin{lstlisting}[language=Python]
import statsmodels.api as sm
\end{lstlisting}

\section{Extracting the Data}

We extract columns from the DataFrame:

\begin{lstlisting}[language=Python]
y = df["y"]
X = df["x"]
\end{lstlisting}

Notice:

\begin{itemize}
    \item uppercase \(X\) is used for independent variables,
    \item lowercase \(y\) is used for dependent variables.
\end{itemize}

\section{Adding a Constant Term}

We now call:

\begin{lstlisting}[language=Python]
X = sm.add_constant(X)
\end{lstlisting}

This adds a column of ones.

Why?

Because the model:

\[
y = \beta_0 + \beta_1 x
\]

contains a constant term.

Without this step, the fit would assume:

\[
y = \beta_1 x
\]

which forces the line through the origin.

\section{Creating the Model}

We construct the model using Ordinary Least Squares (OLS):

\begin{lstlisting}[language=Python]
model = sm.OLS(y, X).fit()
\end{lstlisting}

Important:

\begin{itemize}
    \item the dependent variable \(y\) comes first,
    \item the independent-variable matrix \(X\) comes second.
\end{itemize}

\section{What is OLS?}

OLS stands for:

\[
\text{Ordinary Least Squares}
\]

This is the standard linear-regression technique.

The idea is to choose parameters that minimize:

\[
\sum_i (y_i - y_{\mathrm{fit},i})^2
\]

the sum of squared residuals.

\section{Viewing the Results}

The most useful step is:

\begin{lstlisting}[language=Python]
print(model.summary())
\end{lstlisting}

This produces a detailed statistical summary including:

\begin{itemize}
    \item best-fit parameters,
    \item uncertainties,
    \item confidence intervals,
    \item goodness-of-fit measures,
    \item statistical tests.
\end{itemize}

\section{Advantages of StatsModels}

Compared with simpler fitting approaches, \texttt{statsmodels} provides:

\begin{itemize}
    \item cleaner statistical output,
    \item automatic uncertainty calculations,
    \item convenient regression tools,
    \item integration with pandas.
\end{itemize}

\section{Complete Example}

\begin{lstlisting}[language=Python]
import pandas as pd
import statsmodels.api as sm

df = pd.DataFrame({

    "x": [5, 10, 15, 20],
    "y": [5, 12, 20, 27]

})

y = df["y"]

X = df["x"]

X = sm.add_constant(X)

model = sm.OLS(y, X).fit()

print(model.summary())
\end{lstlisting}

\section{Key Takeaways}

\begin{itemize}
    \item Pandas DataFrames organize data using labeled columns.
    \item StatsModels provides advanced regression tools.
    \item \texttt{sm.add\_constant()} adds the intercept term.
    \item OLS performs standard least-squares regression.
    \item \texttt{model.summary()} gives detailed statistical information.
\end{itemize}

\end{document}
